\section{Stereographic QSVT Protocol}
\label{sec:qsvt}

With the block encoding in place, we now show how the
single-qubit QSP protocol lifts to a multi-qubit eigenvalue
transformation.  The key observations are that the phase gates
act only on the ancilla (and thus commute with the system
register), and that the signal operator decomposes as a direct
sum over eigenspaces.

\subsection{Lifting QSP to Multiple Eigenspaces}
\label{subsec:lifting}

For a Hamiltonian with known diagonalization
$H = V^\dagger D V$ and spectral decomposition
$H = \sum_j \lambda_j \ket{j}\bra{j}$, the stereographic
signal operator on the full Hilbert space
$\Hil_a \otimes \Hil_s$ is
\begin{equation}
	\mathcal{S} = \sum_j S_{\lambda_j} \otimes \ket{j}\bra{j},
	\label{eq:multi-signal}
\end{equation}
where $S_{\lambda_j} = U_{\lambda_j}\PauliZ$ is the
single-qubit signal operator for eigenvalue $\lambda_j$.

\begin{proposition}[Direct-Sum Decomposition]
	\label{prop:direct-sum}
	The phase gate
	$\exp[i\phi\PauliZ] \otimes \Id_s$ commutes with the
	eigenspace projectors $\Id_a \otimes \ket{j}\bra{j}$.
	Consequently, the degree-$d$ stereographic QSVT product
	\begin{equation}
		\mathcal{W}_{\bm{\phi}}
		= \prod_{k=0}^{d}
		\bigl(\exp[i\phi_k\PauliZ] \otimes \Id_s\bigr)
		\cdot \mathcal{S}
		\label{eq:qsvt-product}
	\end{equation}
	decomposes as
	\begin{equation}
		\mathcal{W}_{\bm{\phi}}
		= \sum_j W_{\bm{\phi}}(\lambda_j) \otimes
		\ket{j}\bra{j},
		\label{eq:qsvt-decomposition}
	\end{equation}
	where $W_{\bm{\phi}}(\lambda_j)$ is the single-qubit
	stereographic QSP product evaluated at $r = \lambda_j$.
\end{proposition}

\begin{proof}
	Since $\exp[i\phi\PauliZ] \otimes \Id_s$ is diagonal in
	$\Hil_a$ and acts as the identity on $\Hil_s$, it commutes
	with $\Id_a \otimes \ket{j}\bra{j}$.  Therefore each
	alternating factor in the QSVT product can be grouped within
	the eigenspace projection, giving the direct-sum
	decomposition.
\end{proof}

\subsection{The Full Protocol Circuit}
\label{subsec:full-circuit}

For a Hamiltonian with known diagonalization
$H = V^\dagger D V$, the full stereographic QSVT circuit is:
\begin{equation}
	\begin{quantikz}[column sep=0.2cm, row sep=0.4cm]
		\lstick{\footnotesize$\ket{0}_a$} & \gate{\phi_d} & \gate{W_j} & \ \cdots\ \qw & \gate{W_j} & \gate{\phi_0} & \meter{} \\
		\lstick{\footnotesize$\ket{\psi}_s$} & \gate[2]{V} & \ctrl{-1}\vqw{1} & \ \cdots\ \qw & \ctrl{-1}\vqw{1} & \gate[2]{V^\dagger} & \qw \\
		& & \ctrl{-1} & \ \cdots\ \qw & \ctrl{-1} & & \qw
	\end{quantikz}
	\label{eq:full-protocol-circuit}
\end{equation}
Each $\phi_k$ box denotes $R_z(\phi_k) = \exp[i\phi_k\PauliZ]$
on the ancilla, and $W_j \equiv W(\tilde{\lambda}_j)$ is the
stereographic signal unitary conditioned on eigenstate $\ket{j}$.

\begin{theorem}[Eigenvalue Transformation]
	\label{thm:eigenvalue-transform}
	Let $H = \sum_j \lambda_j \ket{j}\bra{j}$ be a Hermitian
	operator with stereographic block encoding~$U$, and let
	$\bm{\phi}$ be a set of QSP phases realizing polynomials
	$(P, Q)$ with $|P|^2 + |Q|^2 = 1$.  Then the stereographic
	QSVT protocol produces the state
	\begin{equation}
		\mathcal{W}_{\bm{\phi}}\ket{0}_a\ket{j}_s
		= \bigl(P(\tilde{\lambda}_j)\ket{0}
		+ Q(\tilde{\lambda}_j)\ket{1}\bigr)\ket{j}_s,
		\label{eq:eigenvalue-state}
	\end{equation}
	where $\tilde{\lambda}_j = \lambda_j/\sqrt{1+\lambda_j^2}$.
	Pauli-measurement decoding of the ancilla yields
	\begin{equation}
		f(\lambda_j)
		= \frac{P(\tilde{\lambda}_j)}{Q(\tilde{\lambda}_j)},
		\label{eq:eigenvalue-transform}
	\end{equation}
	a rational function of $\lambda_j$ of total degree
	$\leq d$, with circuit depth $\mathcal{O}(d)$ independent
	of $\|H\|$.
\end{theorem}

\begin{proof}
	By Proposition~\ref{prop:direct-sum}, the QSVT product
	decomposes into independent single-qubit QSP products in
	each eigenspace.  The result then follows from the
	single-qubit theory~\cite{KharaziLeytonOrtega2026a}.
\end{proof}

\subsection{Depth Advantage}
\label{subsec:depth-advantage}

\begin{corollary}[Depth Comparison]
	\label{cor:depth}
	For a degree-$d$ polynomial transformation of the
	eigenvalues of a Hermitian operator $H$ with
	$\|H\| = \alpha$:
	\begin{enumerate}
		\item Standard QSVT requires circuit depth
		      $\mathcal{O}(d \cdot \alpha)$.
		\item Stereographic QSVT requires circuit depth
		      $\mathcal{O}(d + C_V)$, where $C_V$ is the
		      gate complexity of the diagonalization $V$.
	\end{enumerate}
	The depth saving is a multiplicative factor of~$\alpha$
	in the QSP component, at the cost of requiring the
	diagonalization circuit~$V$.
\end{corollary}
